% ================================= CONCLUSAO =================================
\section{Conclusão}\label{sec:conclusao}

Este estudo mediu o valor incremental de \textit{random forest} e XGBoost sobre uma regressão logística com as cinco variáveis do Escore de Mayo na predição de óbito em até cinco anos na cirrose biliar primária. Num desenho 2 x 2, com as mesmas partições e o mesmo orçamento de ajuste para todos os modelos, nenhum modelo de AM demonstrou ganho, nem com as 5 variáveis nem com as 17, e a referência teve AUROC nominalmente próxima da dos demais modelos, sem diferença demonstrada, e estimativas de calibração próximas das ideais. Neste estudo de caso, um modelo de regressão com cinco variáveis de rotina vale mais a pena: o ganho do AM não foi demonstrado, e seus custos, em transparência e, com 17 variáveis, em coleta de dados, são concretos. O braço de controle mostrou ainda que o enquadramento da competição do Kaggle, com o tempo de seguimento como entrada, infla o desempenho com uma informação indisponível no momento da predição.

A continuidade natural é tratar o problema como análise de sobrevivência, comparando o modelo de Cox com \textit{random survival forests} sem excluir os pacientes censurados, e repetir a comparação em coortes atuais de pacientes tratados com ácido ursodesoxicólico, em que escores como o GLOBE e o UK-PBC são as referências clínicas. Em ambos os casos, a pergunta deve continuar sendo a do valor incremental: quanto o novo modelo acrescenta ao que a clínica já usa, e a que custo.
